% ============================================================
% Abstract
% ============================================================
\begin{abstract}
Bayesian optimization (BO) is widely used for expensive black-box
functions, where a Gaussian process (GP) surrogate drives an acquisition
function that balances exploration of uncertain regions against
exploitation of promising ones.
A fundamental open question is whether the GP's uncertainty estimate
alone—without any exploitation signal—is sufficient to outperform blind
random search.

I address this question with a systematic benchmark of four acquisition
strategies: Expected Improvement (EI), Upper Confidence Bound (UCB),
Uncertainty (maximum posterior standard deviation), and Random.
All four strategies share the same normalized GP surrogate and are
evaluated on six problems spanning dimensions 2 to 10, mixing classical
synthetic functions with physics-based engineering test functions.
Each strategy–problem pair is repeated over 10 independent seeds,
yielding 240 BO runs total.
A Friedman–Nemenyi non-parametric test ($N = 60$ blocks, $k = 4$)
provides statistically robust comparisons.

The Friedman test strongly rejects equal performance
($\chi^2(3) = 128.82$, $p = 9.73 \times 10^{-28}$).
Nemenyi post-hoc comparisons reveal two statistically separated
equivalence classes: balanced strategies $\{\text{EI, UCB}\}$
significantly outperform pure-exploration strategies
$\{\text{Uncertainty, Random}\}$ (all cross-class $p < 10^{-6}$),
while Uncertainty and Random are indistinguishable from each other
($p = 0.974$).
This Uncertainty–Random tie demonstrates that exploitation of the GP's
posterior mean—not the uncertainty estimate alone—is the necessary
condition for sample-efficient BO.
A secondary finding is that GP input normalization to known problem
bounds is a first-class engineering requirement on multi-scale problems:
omitting it degrades BO to random-search performance on Piston (7D),
while leaving the Random baseline unaffected, isolating the fault to
the surrogate.
\end{abstract}
